% l4_contracts.tex --- Layer 4 (Output Contracts), P4 draft, 2026-10-17.
% Subsections for \section{Layer 4 --- Output Contracts: Validity Is Not Truth}\label{sec:contracts}.
% Placeholder convention (same as sibling sections): \cite{p2[<arxiv-id>]} for IDs verified in the
% P2 sweep, wired to named keys at merge. p2 IDs cited here: 2305.13971, 2609.33887, 2609.33843,
% 2609.34024, 2610.18529, 2610.19811, 2610.02076, 2610.02267. Already wired: jsonschemabench2025,
% xgrammar2024, rauk2026jevbenchmark. Vendor documentation cited in footnotes per the L5 precedent.
% NOTE: 2609.33843 carries an author-similarity flag in the Laya audit (item f) --- cite as one
% reproduction, not consensus; miscalibration direction reported contested per audit item (a).

\subsection{The contract is the interface the software can rely on}
\label{sec:l4-contract}

Layer~4 asks what a component \emph{returns} to the software around it. We call the declared,
externally checkable shape of that return an \emph{output contract}, ordered in four contract
tiers (T1--T4), numbered as in \cref{fig:layer4-contract}.
\emph{Free-form text} is the weakest: downstream code can act on it only through another parser or
a model, so the contract is ``whatever the consumer can extract'' (T1). \emph{Schema-constrained output}
adds a grammar or JSON schema the return must satisfy (T2). \emph{Typed decisions} remove text entirely:
the component returns a distribution over a fixed set of typed choices (a probability per option
or per stated claim), so the decision score is what thresholds consume (T3).
\emph{Structured metadata} --- citations, spans, tool-call records, confidence fields --- layers
evidence onto any of the above (T4). The tier fixes what later layers can check and select over
(\cref{sec:selection}): a validator can test a grammar, a threshold can consume a score, and
neither can test a sentence. That is the handoff promised in \cref{sec:l2-handoff}, where representations
decide what the system computes over and contracts decide what can be verified about what it
returns.

\subsection{The validity half: constrained decoding, and what validity measures}
\label{sec:l4-validity}

Grammar-constrained decoding makes the validity half constructive: at each step it suppresses every
token that would violate the grammar. \evindep{Constrained decoding was introduced to close a specific gap: language models fail to produce complex structures reliably when not
finetuned to the required format, and masking invalid continuations restores structural correctness
without retraining \cite{id230513971}.} Serving engines followed: XGrammar constrains generation
against grammars and JSON schemas at practical overhead \cite{xgrammar2024}. The guarantee is
precise and narrow: whatever the engine emits parses under the declared schema. What it is worth
is measured. \evindep{JSONSchemaBench separates coverage (which schema features an engine
actually supports), output quality, and generation efficiency, rather than treating ``supports
JSON'' as one property \cite{jsonschemabench2025}.} The separation matters: a schema an engine
cannot fully express degrades silently to a weaker contract, and latency can make a valid output
unusable. Our first load-bearing conclusion, by synthesis rather than citation: \evsynth{schema compliance is table stakes: validity is not correctness. Whether the values filling the schema are true is a distinct property
that no decoding-time guarantee addresses, because the grammar constrains form and says nothing
about content.} A perfectly formed JSON object can be confidently wrong in every field.

\subsection{The truth half: typed decisions and calibrated probabilities}
\label{sec:l4-typeddecisions}

Decision-oriented models occupy the typed-decision tier by construction. \evindep{The
Single-Shot/System-One family returns no text: given a state and typed questions, the model returns
a choice from fixed options, a position on a rubric, or the probability that a statement is true,
and independent benchmarking evaluates these returns on held-out data
\cite{rauk2026jevbenchmark}.} The vendor architecture documents the interface in its own vocabulary
--- typed \texttt{choice}, \texttt{score} and \texttt{noul} primitives, \texttt{noul} (the
vendor's spelling) being a
$P(\text{true})$ query --- and reports an RLCD recipe built on proper scoring rules (log score,
spherical score, ranked probability score) whose object is a truthfully reported
distribution.\footnote{\texttt{github.com/NandhaKishorM/laya}
documentation, accessed 2026-10-17; evidence class \evvendor{}.} The contract does not remove the
truth problem; it relocates it, from ``is the text true?'' to ``are the probabilities honest?''
That second question is measurable. \evindep{Independent evaluations now span decision accuracy and
calibration in medicine \cite{id260934024,id261018529}, binary decision scores under a risk
budget \cite{id261019811}, and fine-tuning comparisons \cite{id261002076}.} Vendor
probability claims enter as hypotheses, not facts; that is the evidence discipline of this survey.

The discipline has already caught a live disagreement. \evvendor{Vendor documentation warns
that shipped checkpoints are over-confident and need temperature refitting before gating relies on
them.} \evindep{One reproduction study measures the shipped typed-decisions checkpoint as
uniformly \emph{under}-confident, signed confidence--accuracy gap $-0.214$, improved by a single
fitted sharpening temperature \cite{id260933843}.} The direction of miscalibration is therefore
reported as \emph{contested} (an open ledger entry in \cref{sec:protocol}, cited per checkpoint
and never as a family-level fact), and the sign decides which way a confidence-gated cascade
fails: an over-confident model passes wrong answers through the gate; an under-confident one
abstains on correct ones. A gate designed on the wrong sign misprices both its accepted-error
budget and its coverage.

\subsection{Thresholds are policy; option sets set the scale}
\label{sec:l4-thresholds}

Typed probabilities become gating machinery only through a threshold, and the operational
literature has converged on the stance this paper adopts: ``a threshold is a policy you choose
from measured accuracy at that coverage on your data, not a property of the model''\footnote{Ibid.,
``Automated Confidence Gating,'' which in our framework is derived from the risk budget of
\cref{sec:l5-calibrated}.} Two constraint classes --- option count and option order --- enter the
ledger as separately sourced rows; no single law covers them. \evvendor{The vendor's decision-model
card advises keeping \texttt{choice} questions under roughly twenty options within a fixed
per-label budget, with sharp falloff beyond; the mechanism is budget trimming, so we record
guidance, not a measured cliff.} \evvendor{On a 77-label intent benchmark the vendor reports 0.425
against 0.870 intent accuracy for the Jev-class model. The comparison is uncontrolled and cross-source; the
vendor itself flags it as indicative.} \evindep{Independently, paired evaluation finds
the same family changes 30\% of its answers under option-order reversal and reaches 31\% accuracy
at fifty nearest-neighbour tools \cite{id261002267}.} Option count and option order are
therefore contract-design variables. Any accuracy or threshold claim is only meaningful at the
option configuration it was measured on.

\subsection{The hub: risk-controlled serving as the L4--L5 bridge}
\label{sec:l4-hub}

The contract tier creates the preconditions for risk-controlled operation: typed, checkable
outputs with a declared correctness measure are what calibration needs. The cleanest current
construction is risk-budgeted configuration selection (Redline). \evtheory{Given a reference
configuration and a candidate, a finite-sample procedure selects operating points --- acceptance
thresholds, buffer depth, schedule, checkpoint, precision --- such that the
\emph{reference-relative joint risk}, the probability over prompts that the reference answers
correctly while the candidate does not, stays within a user-chosen budget with high probability,
with no distributional assumptions \cite{id260933887}.} Its scope is as sharp as its strength:
relative to a reference and a chosen correctness measure, and by construction \emph{not} a
factual-correctness guarantee; speed is bought at a bounded disagreement rate, not truth. The
budget is also task-dependent in the authors' measurements (math clears at a tighter budget
than code), which is why we treat it as an L0 input rather than a constant. Full treatment is
\cref{sec:l5-bridge}. The contract-side lesson is that a contract fixes what ``accurate'' means
for the return, which is what lets ``faster but less accurate'' become a number an operator can
budget.

\subsection{What Layer 4 does not guarantee}
\label{sec:l4-limits}

Each contract tier constrains \emph{form}, and only form. Free-form text guarantees nothing
checkable; schema-constrained output guarantees parseability, not that the parse is true; typed
decision probabilities make confidence auditable, not accurate; structured metadata makes evidence
\emph{attachable}, not attached or faithful. \evsynth{An NCM-class staged tokenizer (\cite{kumbong2026ncm}) sharpens the same split one stage
earlier: the token-exactness vs.\
meaning-equivalence boundary is a contract statement to define explicitly --- reconstruction
fidelity on the code side is not meaning fidelity on the read
side.} \evsynth{A schedule change also relocates the error surface: grammar and schema gates
operate on decoded tokens and never see the latent codes, so a substitution becomes observable
only after decoding, and the latent channel remains unaudited while acceptance machinery must be
anchored at the decoded surface.} \evaurep{The authors' conditioning-corruption training --- their
code-error robustness measure --- is the corresponding hardening move: robustness bought inside
the representation rather than enforced at the contract.} Semantics enter only through Layer~5's checks ---
grounding, validators, calibrated selective prediction, escalation to human review (\cref{sec:assurance}) --- which are
themselves fallible components, as the assurance layer's own failure envelope concedes
(\cref{sec:l5-envelope}). The asymmetry is structural: form is
checkable by construction; truth is a separate property. It is also why the selection framework
(\cref{sec:selection}) treats the contract tier as a design decision, not a formatting preference.
